\chapter*{Résumé}
\addcontentsline{toc}{chapter}{Résumé}

\begin{otherlanguage}{french}
Les nomenclatures logicielles, ou SBOM, sont de plus en plus utilisées pour
renforcer la transparence de la chaîne d'approvisionnement logicielle et la
gestion des vulnérabilités. Toutefois, une SBOM syntaxiquement valide n'est
pas automatiquement digne de confiance, en particulier lorsqu'elle est
reconstruite à partir d'une application Android compilée, sans accès au code
source ni aux manifestes de dépendances. La compilation, l'optimisation,
l'obfuscation, les alias de familles d'artefacts, les lacunes des bases de
références et la similarité entre versions peuvent affaiblir ou contredire les
preuves qui soutiennent une déclaration de composant.

Ce mémoire étudie la reconstruction fiable de SBOM comme un problème
d'attribution de composants logiciels fondé sur les preuves. Il organise les
travaux existants sur la détection de bibliothèques tierces Android et la
récupération de composants binaires selon des taxonomies de techniques, de
sources de preuves, de menaces de transformation et de pratiques
d'évaluation. À partir des lacunes ainsi mises en évidence, il propose un
cadre conceptuel dans lequel chaque déclaration reconstruite conserve
l'identité du composant, la version ou l'ensemble de versions, la provenance
des preuves, leur qualité, le risque de transformation, le niveau de
confiance, l'ambiguïté, la classe de validation et l'interprétation de
sécurité autorisée en aval.

Un système d'analyse APK-vers-SBOM implémenté sert d'étude de cas délimitée. Ses
preuves contrôlées, propres et partielles validées par les sources, ses bornes
supérieures connues par les sources, ainsi que ses résultats expérimentaux,
indicatifs et limites, montrent pourquoi l'identité de bibliothèque, la
version exacte, la couverture des références, le rappel du moteur de
correspondance et l'autorité des vulnérabilités doivent rester séparés. Le
cadre est évalué par l'étude de cas : il préserve ces catégories de preuves et
empêche toute augmentation d'autorité lorsqu'une modification apparemment
bénéfique échoue aux critères d'ambiguïté, de contrôle ou de promotion. Le
cadre proposé privilégie donc une connaissance partielle auditable et
l'abstention explicite plutôt qu'une certitude non étayée, et définit un
programme de recherche pour la fusion calibrée des preuves, les SBOM sensibles
à la provenance, les bancs d'essai réalistes et la récupération de composants
natifs ou chargés dynamiquement.
\end{otherlanguage}

\vspace{0.5cm}
\noindent\textbf{Mots-clés ---}
\begin{otherlanguage}{french}
nomenclature logicielle, chaîne d'approvisionnement logicielle, Android,
attribution de composants, provenance, fusion de preuves, confiance,
incertitude, gestion des vulnérabilités.
\end{otherlanguage}
